\scanpage{7}

vaut mieux dès lors exiger que ces \unsure{isos|arêtes} soient \uline{homographiques} par morceaux....

Dans le premier cas (exigeant que les disques de la découpe soient poncturés...) on trouve, en supposant $\pi_0(K_0)\longrightarrow\pi_0(X_0)$ surjectif, \deleted{et quitte \unsure{à éliminer} des bords isolés....}
\[
\text{nb de modules \inserted{réels}}\quad N=\left(\sum_{\substack{f\in\pi_0(X\backslash K)=F\\f\text{ pas \uline{un} disque}}}
(6g(f)+4\nu(f)\deleted{\unsure{+}}-6)\right)\deleted{+\unsure{\mu}}
\]
\deleted{$\mu$ étant le nb de ``faces'' $f$ qui sont des \deleted{disques}\inserted{\unsure{anneaux|couronnes}}.}

$\checkmark/$ si $g\ge 2$

\margin{On suppose $X_0^0$ connexe ainsi que $K$ aussi. $\lambda$ la connexité de $K$}
\[
\begin{aligned}
3g-3\qquad\text{où}\quad 2-2g&=\chi=\sum_{f\in F}\bigl(2-2g(f)-\nu(f)\bigr)\\
&\hspace{8em}+1-\lambda\\
\text{i.e.}\qquad 2g-2&=\sum_{f\in F}\bigl(2g(f)-2+\nu(f)\bigr)\\
&\hspace{8em}+(\lambda-1)\\
g-1&=\tfrac12\sum_{f\in F}\bigl[2g(f)-2+\nu(f)-\bigr]\\
&\hspace{8em}+\tfrac{\unsure{3|1}}2(\lambda-1)\\
3(g-1)&=\tfrac32\sum_{f\in F}\bigl[2g(f)-2+\nu(f)\bigr].\\
&\hspace{8em}+\tfrac32(\lambda-1)\\
&=\sum_{f\in F}\bigl[3(g(f)-1)+\tfrac32\deleted{\unsure{2}}\nu(f)\bigr]\\
&\hspace{8em}+\tfrac32(\lambda-1)
\end{aligned}
\]
\[
(g\ge2)\qquad N-(3g-3)=\sum_{\substack{f\in F\\f\text{ pas un disque}}}
\bigl(3g(f)+\tfrac52\nu(f)-3\bigr)\deleted{\unsure{-}}+
\deleted{\sum_{\substack{f\in F\\f\ \mathrm{disque}}}\unsure{1}}
\tfrac32(\mu+\unsure{1-\lambda})\overset{?}{\ge}0
\]
$\mu$ étant le nb de $f$ qui sont des disques
\[
(g=1)\qquad N-1=\sum_{\substack{f\ \text{pas un}\\\text{disque}\\\text{[NB }g(f)=0\ \text{tjrs]}}}
(\underbrace{4\nu(f)-6}_{\geqslant 2})-1
\]
% archive sheet number: 6
